% Source: direct-forced-construction-20260910.md, section 8, equations (37)--(43).
% Pressure normalization: section 2, equation (5), of that source.
\subsection{From spatial-height integrals to the forced mixed intersection}
\label{r:sec:forced-domain-propagation}

Let \((u,p)\) be the classical forced history on every strict interval
\([0,S]\), \(S<T<\infty\), with fixed \(\nu>0\),
\(u_0\in H^\infty_\sigma(\R^3)\), and the prescribed smooth force whose
mixed derivatives decay rapidly in space uniformly through \([0,T]\).
Assume explicitly that every zeroth-row completed height has its
whole-interval integral:
\begin{equation}
 \int_0^T C^f_{0,k}(t)\dd t
 =\frac12\int_0^T\norm{\Lambda^{k+1/2}u(t)}_2^2\dd t<\infty,
 \qquad k=0,1,2,\ldots .
 \label{r:eq:forced-spatial-domain-hypothesis}
\end{equation}
The forced momentum equation carries this spatial information to the
temporal rows. The first step is to recover the low frequencies omitted
by the homogeneous heights. Put
\[
 K_T=\norm{u_0}_2+\int_0^T\norm{f(t)}_2\dd t,
 \qquad \mathsf S_m=\norm{u}_{L^2(0,T;H^m)}.
\]
The kinetic bound gives \(\sup_{t<T}\norm{u(t)}_2\le K_T\).
For each integer \(m\ge0\), split the Fourier integral at \(|\xi|=1\).
Below that radius use the kinetic bound; above it use
\((1+|\xi|^2)^m\le2^m|\xi|^{2m+1}\). This gives
\begin{equation}
 \mathsf S_m^2
 \le2^mTK_T^2+2^{m+1}\int_0^T C^f_{0,m}(t)\dd t<\infty.
 \label{r:eq:forced-domain-spatial-integrals}
\end{equation}

These time integrals bound the total change in each spatial norm.
Indeed, the projected equation is
\[
 u_t=\nu\Delta u-\Pp[(u\cdot\nabla)u]+\Pp f.
\]
The transport estimate \eqref{r:eq:transport-product} places its nonlinear
term in \(L^1(0,T;H^m)\), with norm at most
\(C_m\mathsf S_{m+2}^2\). Its viscous term has \(L^1H^m\) norm at most
\(\nu\sqrt T\,\mathsf S_{m+2}\). Integrating the equation from the
initial datum now yields
\begin{equation}
 \begin{aligned}
 \sup_{t<T}\norm{u(t)}_{H^m}
 &\le\mathsf U_{0,m},\\
 \mathsf U_{0,m}
 &:=\norm{u_0}_{H^m}+\nu\sqrt T\,\mathsf S_{m+2}
       +C_m\mathsf S_{m+2}^2
       +\norm{\Pp f}_{L^1(0,T;H^m)}<\infty.
 \end{aligned}
 \label{r:eq:forced-domain-base-row}
\end{equation}
The integral of this \(L^1H^m\) right side supplies the continuous
representative, so the bound does not assume a terminal trace.

Every finite spatial order is now available in the base row.
Differentiating the forced equation retains every transport product:
\begin{equation}
 \begin{aligned}
 V_{n+1}&=\nu\Delta V_n-\Pp B_n+\Pp F_n,\\
 B_n&=\sum_{a=0}^n\binom na(V_a\cdot\nabla)V_{n-a},
 \qquad V_n=\partial_t^nu,\quad F_n=\partial_t^nf.
 \end{aligned}
 \label{r:eq:forced-domain-row-recurrence}
\end{equation}
Suppose that rows \(0,\ldots,n\) have uniform bounds at every finite
spatial order. Applying \eqref{r:eq:transport-product} term by term in
this finite sum gives the next constants:
\begin{equation}
 \begin{aligned}
 \mathsf U_{n+1,m}:={}&\nu\mathsf U_{n,m+2}\\
 &+C_m\sum_{a=0}^n\binom na
       \mathsf U_{a,m+2}\mathsf U_{n-a,m+2}
       +\norm{\Pp F_n}_{L^\infty(0,T;H^m)}.
 \end{aligned}
 \label{r:eq:forced-domain-row-bounds}
\end{equation}
The prescribed force makes the last term finite. Induction on \(n\)
proves
\(\sup_{t<T}\norm{V_n(t)}_{H^m}\le\mathsf U_{n,m}<\infty\)
for every fixed pair \(n,m\). At each step only finitely many higher
spatial orders occur, and each of them was already available in the
earlier rows. Integrating these uniform bounds gives
\begin{equation}
 \sum_{n=0}^r\norm{V_n}_{L^2(0,T;H^m)}^2
 \le T\sum_{n=0}^r\mathsf U_{n,m}^2<\infty,
 \qquad
 u\in\mathscr I_T=\bigcap_{r,m\ge0}H^r(0,T;H^m).
 \label{r:eq:forced-domain-mixed-intersection}
\end{equation}
Conversely, choosing spatial order \(k+1\) in this intersection gives
\eqref{r:eq:forced-spatial-domain-hypothesis}. The full spatial-height
domain and the mixed intersection are equivalent for this forced
history.

Pressure completes the momentum rows without changing the prescribed
force. Write \(Q_n=\partial_t^np\), and use the decaying pressure
normalization
\begin{equation}
 \begin{aligned}
 \phi_n&=-(-\Delta)^{-1}\nabla\cdot F_n,\\
 Q_n&=R_iR_j\sum_{a=0}^n\binom na
                     V_{a,i}V_{n-a,j}+\phi_n,\\
 \nabla Q_n&=(I-\Pp)(F_n-B_n),
 \end{aligned}
 \label{r:eq:forced-domain-pressure}
\end{equation}
where \(R_i\) are the Riesz transforms and repeated spatial indices
are summed. The force potential is finite in every \(H^M\).
At low frequencies, \(\widehat F_n\) is uniformly bounded and its
inverse-derivative multiplier is controlled by
\(\int_{|\xi|\le1}|\xi|^{-2}\dd\xi=4\pi\).
At high frequencies, the prescribed Sobolev norms suffice. The smooth
time dependence of the source gives
\(\phi_n\in C([0,T];H^M)\) for every fixed \(n,M\).

The Riesz transforms are bounded on \(H^M\), and
\(H^{M+1}\cdot H^{M+1}\subset H^M\) in three dimensions.
For every integer \(M\ge0\), these product bounds and
Cauchy--Schwarz in time give
\begin{equation}
 \begin{aligned}
 &\norm{Q_n}_{L^1H^M}
       +\norm{\nabla Q_n}_{L^1H^{M-1}}\\
 &\quad\le c_M\sum_{a=0}^n\binom na
       \norm{V_a}_{L^2H^{M+1}}\norm{V_{n-a}}_{L^2H^{M+1}}\\
 &\qquad\quad+\norm{\phi_n}_{L^1H^M}
       +\norm{\nabla\phi_n}_{L^1H^{M-1}}<\infty.
 \end{aligned}
 \label{r:eq:forced-domain-pressure-integrals}
\end{equation}
Here all time norms are over \((0,T)\). In particular, \(M=0\)
uses \(H^{-1}\) for the gradient norm.
The velocity part of each finite rectangle satisfies
\begin{equation}
 R_{N,M}
 \le T\sum_{n=0}^N
 \left(\mathsf U_{n,M+1}^2+
       \mathsf U_{n+1,\max\{M-1,0\}}^2\right)<\infty.
 \label{r:eq:forced-domain-rectangles}
\end{equation}
Each prescribed \(F_n\) belongs to \(L^2(0,T;H^{M-1})\).
Together with \(Q_0,\ldots,Q_N\) in the spaces just proved and
\(F_0,\ldots,F_N\), this retains the full forced rectangle
\((V_0,\ldots,V_{N+1};Q_0,\ldots,Q_N;F_0,\ldots,F_N)\).
Every retained row obeys
\[
 V_{n+1}+B_n+\nabla Q_n-\nu\Delta V_n=F_n,
 \qquad \nabla\cdot V_n=0.
\]
Restriction to a smaller rectangle preserves all derivatives, pressure
terms, and force coordinates because these entries come from the
given history.

The next temporal row also controls approach to the endpoint.
For \(s,t<T\), integration of \(\partial_tV_n=V_{n+1}\) gives
\begin{equation}
 \norm{V_n(t)-V_n(s)}_{H^m}
 \le |t-s|\,\mathsf U_{n+1,m}.
 \label{r:eq:forced-domain-endpoint}
\end{equation}
Each row has an \(H^m\) limit at \(T\). If a trace is taken first in
a higher Sobolev space, its restriction is this limit by uniqueness
of the \(H^m\) limit. The traces assemble into one \(H^\infty\) jet,
with \(u_*=V_0(T)\in H^\infty_\sigma\).
Product continuity and \eqref{r:eq:forced-domain-pressure} give the
pressure traces and preserve the complete terminal equation:
\[
 V_{n+1}(T)+B_n(T)+\nabla Q_n(T)-\nu\Delta V_n(T)=F_n(T).
\]
Choosing \(m>k+3/2\) in \eqref{r:eq:forced-domain-endpoint} gives uniform
convergence of every spatial derivative of order at most \(k\).
The finite products in the pressure formula give the corresponding
limits for pressure and its gradient.

If the force is prescribed smoothly beyond \(T\), local forced
existence restarts the equation from \(u_*\).
The outgoing initial jet is fixed by
\eqref{r:eq:forced-domain-row-recurrence} and
\eqref{r:eq:forced-domain-pressure}, so it agrees with the incoming
jet. Local uniqueness joins these histories into a smooth solution
through \(T\).
